% problem_id: S047-b21-Q1
% source_id: S047/S048 (numdam/aif3605.pdf)
% source_file: aif3605.pdf
% statement_locator: printed p. journal p. 2325 = PDF page 52 (lower half of the page, immediately after the proof of Theorem A.3 (Wu), which ends on the same page; the definition of M_m that the statement uses is the display two lines above the theorem, also on p. 2325)
% proof_locator: printed p. journal pp. 2325-2326 = PDF pages 52-53 (p. 2325: from 'Proof. --- Similarly to Section 4.1' through the display (A.6); p. 2326: from 'where C := D_{2m+1}/M_m' to the end-of-proof box). The setup identities (A.1)-(A.2) that the proof consumes are on journal p. 2323 = PDF page 50 and are transcribed as marked context.
% representation: PDF; transcribed from rendered page images
% collected_by: carrier rule
% review_status: H2 by 2/2 independent reviewers (the miner claimed H3); 1/2 报不忠实
% status: complete (statement + author proof verbatim + reason + locators)
%
%% ===========================================================================
%% Q1  --  human-proof-corpus candidate transcription
%% ---------------------------------------------------------------------------
%% Source PDF : /disks/sata1/yupeng/human-proof-corpus/source-cache/
%%              registry-sources/numdam/aif3605.pdf   (host Yupeng-orx:443)
%% Paper      : Dmitry Chelkak, Clement Hongler & Remy Mahfouf,
%%              "Magnetization in the zig-zag layered Ising model and
%%               orthogonal polynomials",
%%              Annales de l'Institut Fourier 74, no 6 (2024), p. 2275-2330.
%%              DOI 10.5802/aif.3605
%% Collected  : Theorem A.4  (statement + COMPLETE proof), journal pp. 2325-2326.
%%              Transcribed with it, clearly marked as context:
%%                - the definition of M_m (same page, immediately above the
%%                  theorem),
%%                - the setup identities (A.1) and (A.2) for the full-plane
%%                  observable X_{[v,u]} (journal p. 2323), which the proof
%%                  consumes,
%%                - Remark A.5 (journal p. 2326), the companion new identity
%%                  (A.7), which has NO proof of its own ("It is not hard to
%%                  repeat the proof of Theorem A.4 in this situation").
%% Pages      : journal page = PDF page + 2273.
%%              statement : PDF p. 52 = journal p. 2325
%%              proof     : PDF pp. 52-53 = journal pp. 2325-2326
%%              context   : PDF p. 50 = journal p. 2323 ; PDF p. 53 = journal
%%                          p. 2326 (Remark A.5)
%% Method     : pdftotext used for LOCATING only (its text layer crushes the
%%              symbols).  Everything below is transcribed by hand from page
%%              images rendered with
%%                pdftoppm -f <p> -l <p> -r 150 -png          (whole pages)
%%                pdftoppm -f <p> -l <p> -r 400 -png -x .. -W ..  (crops)
%%              and read as images; the 400 dpi crops were used to settle
%%              every exponent, every boldface vector and every "not in 2Z".
%% Typography : reproduced as printed.  The paper sets vectors in boldface
%%              (\mathbf{v}, \mathbf{u}), the half-plane in blackboard bold
%%              with a diamond (\mathbb{H}^{\diamond}), and uses the curly
%%              inequalities \geqslant / \leqslant.  Copied as printed.
%% Slips      : kept and annotated at the point of occurrence:
%%              (a) journal p. 2326 prints "the discrete Schwartz reflection"
%%                  (the standard name is Schwarz; the same appendix spells it
%%                  "Schwarz" on p. 2323);
%%              (b) journal p. 2326 prints "spinor ... which satisfy the
%%                  boundary conditions" (plural verb for a singular subject);
%%              (c) journal p. 2278 prints the double negative "we were unable
%%                  to find neither Theorem A.4 nor the identity (A.7)".
%% ILLEGIBLE  : none -- every glyph on the transcribed pages was legible.
%% ===========================================================================

\documentclass[11pt]{article}
\usepackage[utf8]{inputenc}
\usepackage[T1]{fontenc}
\usepackage{amsmath,amssymb}
\usepackage{amsthm}
\usepackage[margin=1in]{geometry}
\allowdisplaybreaks

\newcommand{\HHd}{\mathbb{H}^{\diamond}}
\newtheorem*{thmA}{Theorem A.4.}

\begin{document}

\begin{center}
{\large\bfseries Q1 --- Theorem A.4 of Chelkak--Hongler--Mahfouf, \emph{Magnetization in
the zig-zag layered Ising model and orthogonal polynomials}}\\[2pt]
{\small Ann.\ Inst.\ Fourier \textbf{74}, no 6 (2024), p.\ 2275--2330; DOI 10.5802/aif.3605;
source \texttt{aif3605.pdf}}\\[6pt]
{\small Verbatim transcription. \textsf{[CONTEXT]} blocks are not part of the collected
unit; they are printed so that the unit can be read on its own.}
\end{center}

\medskip
\hrule
\medskip

%% ===========================================================================
%% CONTEXT (journal p. 2323 = PDF p. 50)
%% ===========================================================================

%% --- page 2323 (PDF p. 50) ---

\textsf{[CONTEXT]} The sentences that introduce the appendix's new identities, and the
setup of the full-plane observable $X_{[\mathbf{v},\mathbf{u}]}$ with its values
(A.1) and its Laplacian at the branch points (A.2) --- both of which the proof of
Theorem A.4 uses:

To obtain the latter, we use a simple Schwarz reflection argument instead
of applying our main result, Theorem 1.1, in the spirit of Section 4.3. This
gives a set of exact identities (see Theorem A.4 and Remark A.5) between
$M_m$ and diagonal correlations in the full plane which appear to be new.

Let $n\in\mathbb{N}_0$ and assume that the $\frac{\pi}{4}$-rotated square grid is shifted
so that its vertices (resp., centers of faces) form the lattice $(-n-\frac12+k,s)$
(resp., $(n+\frac12+k,s)$) with $k,s\in\mathbb{Z}$ and $k+s\in 2\mathbb{Z}$. Let
\[
  D_n := \mathbb{E}[\sigma_{(-n+\frac12,0)}\sigma_{(n+\frac12,0)}]
\]
be the (infinite-volume limit of the) diagonal spin-spin correlation at distance of
$n$ diagonal steps. Denote $\mathbf{v}:=(-n-\frac12,0)$, $\mathbf{u}:=(n+\frac12,0)$ and let
\[
  V(k,s) := X_{[\mathbf{v},\mathbf{u}]}((k,s)),\qquad
  k,s\in\mathbb{Z},\ k+s+n\in 2\mathbb{Z},
\]
recall that $V$ is a spinor on the double covers branching over $\mathbf{v}$ and
$\mathbf{u}$. It follows from Proposition 2.2 (or, equivalently, Proposition 2.4) that
$V$ satisfies the standard discrete harmonicity condition $[\Delta V](k,s)=0$ for all
$k,s$ except at the points $(\pm n,0)$ near the branchings, where
\[
  \relax[\Delta V](k,s) := -V(k,s) + \frac14\big[V(k-1,s-1)+V(k+1,s-1)
                                        + V(k-1,s+1)+V(k+1,s+1)\big].
\]
It directly follows from the definition of the observable $X_{[\mathbf{v},\mathbf{u}]}$
and the self-duality of the critical model that
\begin{equation}
  \tag{A.1}
  V(-n,0) = V(n,0) = D_n\,.
\end{equation}
Moreover, a straightforward computation similar to the proof of Proposition 2.2 implies
that
\begin{equation}
  \tag{A.2}
  \begin{aligned}
    \relax[\Delta V](\pm n,0) &= -\frac12 D_{n+1} &&\text{if } n\geqslant 1\,,\\
    \relax[\Delta V](0,0) &= -D_1 &&\text{if } n=0\,.
  \end{aligned}
\end{equation}

%% ===========================================================================
%% THE COLLECTED UNIT (journal pp. 2325-2326 = PDF pp. 52-53)
%% ===========================================================================

%% --- page 2325 (PDF p. 52) ---

We now move on to an explicit expression for the magnetization in the $(2m)$-th column
of the zig-zag half-plane $\HHd$ with ``$+$'' boundary conditions:
\[
  M_m := \mathbb{E}^{+}_{\HHd}[\sigma_{(-2m-\frac12,0)}]\,.
\]

\begin{thmA}
--- The following identities are fulfilled for all $m\in\mathbb{N}_0$:
\begin{equation}
  \tag{A.5}
  \frac{M_{m+1}}{M_m} = \frac{D_{2m+2}}{D_{2m+1}}\,,\qquad
  M_m = \left(\frac{2}{\pi}\right)^{m}\cdot
        \prod_{k=1}^{2m-1}\left(1-\frac{1}{4k^2}\right)^{\lfloor\frac{k}{2}\rfloor-m}.
\end{equation}
\end{thmA}

\begin{proof}
--- Similarly to Section 4.1, let $\mathbf{v}=(-2m-\frac32,0)$ and
\[
  H(-k,s) := X_{[\mathbf{v}]}((-k,s)),\qquad
  k\in\mathbb{N}_0,\ s\in\mathbb{Z},\ k+s\notin 2\mathbb{Z}
\]
be the half-plane fermionic observable. This is a bounded discrete harmonic
(except at $(-2m-1,0)$) spinor on the double cover of $\HHd$ branching over
$\mathbf{v}$ which satisfy the boundary conditions
%
%% [sic] the source prints the plural verb "satisfy" with the singular subject
%% "spinor"; kept as printed.
\[
  H(0,s) = 2^{-1/2}\cdot[H(-1,s-1)+H(-1,s+1)],\qquad s\notin 2\mathbb{Z},
\]
on the imaginary line (see (4.3) and (4.4)). Denote
\begin{equation}
  \tag{A.6}
  \begin{array}{ll@{\qquad}l}
    V(\pm k,s) := CH(k,s) & \text{if } k\in\mathbb{N}, &
      s\in\mathbb{Z},\ \ k+s\notin 2\mathbb{Z},\\[2pt]
    V(0,s) := 2^{-1/2}\cdot CH(0,s) & \text{if } k=0, &
  \end{array}
\end{equation}
%
%% [layout note] in the source the condition "s in Z, k + s not in 2Z" is set on
%% the right of the two-row (A.6) block, at the height of the row break; the
%% array above reproduces the type block, not the exact vertical offset.

%% --- page 2326 (PDF p. 53) ---

where $C:=D_{2m+1}/M_m$; up to a change of the multiplicative normalization, this is
nothing but the extension of $H$ from the left half-plane to the full plane via the
discrete Schwartz reflection.\footnote{[sic] the source prints ``Schwartz'' here,
while the same appendix prints ``Schwarz'' on journal p.~2323.} By construction, $V$ is
a spinor on the double cover of the full-plane branching over $\mathbf{v}$ and
$\mathbf{u}:=(2m+\frac32,0)$ which is discrete harmonic everywhere (including points on
the imaginary line) except at points $(\pm(2m+1),0)$ near the branchings, where one has
$V(\pm(2m+1),0)=D_{2m+1}$. Therefore, it coincides with the full-plane observable
$X_{[\mathbf{v},\mathbf{u}]}((k,s))$ discussed above. In particular, (A.6) implies the
identity
\[
  \frac12 D_{2m+2} = -[\Delta V](-2m-1,0) = -C\cdot[\Delta H](-2m-1,0)
                  = C\cdot\frac12 M_{m+1}
\]
which is equivalent to the first identity in (A.5). The explicit formula for $M_m$
easily follows from the explicit formula (A.4) by induction.
\end{proof}

%% ===========================================================================
%% CONTEXT (journal p. 2326 = PDF p. 53) -- companion remark, no proof of its own
%% ===========================================================================

\textsf{[CONTEXT]} The companion statement, which the paper announces on p.~2278 as new
together with Theorem A.4 (``Note that we were unable to find neither Theorem A.4 nor
the identity (A.7) in the literature''). It has no proof of its own:

\medskip
\noindent\emph{Remark A.5.} --- Similarly, let $M_{m-\frac12}$ denote the magnetization
in the $(2m-1)$-th column of the critical homogeneous Ising model in the zig-zag plane.
It is not hard to repeat the proof of Theorem A.4 in this situation and to obtain the
identity
\[
  M_{m+\frac12}\,/\,M_{m-\frac12} = D_{2m+1}\,/\,D_{2m},\qquad m\in\mathbb{N}_0,
\]
where we formally set $M_{-\frac12}:=\sqrt2$, this convention is the result of the
additional factor relating the values of the half-plane and the full-plane observables
on the imaginary line via (A.6). By induction, one easily gets the identity
\begin{equation}
  \tag{A.7}
  M_{m+\frac12}M_m = \sqrt2\cdot D_{2m+1},\qquad m\in\frac12\mathbb{N}_0,
\end{equation}
and an explicit formula for $M_{m+\frac12}$, which is similar to (A.5). Finally, a
straightforward analysis gives the asymptotics
\begin{equation}
  \tag{A.8}
  D_n \sim \mathcal{C}_{\sigma}^{2}\cdot(2n)^{-1/4},\qquad
  M_m \sim 2^{1/8}\mathcal{C}_{\sigma}\cdot(2m)^{-1/8},\qquad n,m\to\infty,
\end{equation}
where $\mathcal{C}_{\sigma}=2^{\frac16}e^{\frac32\zeta'(-1)}$ is the same universal
constant as in Theorem 3.9. Note that we prefer to encapsulate the factors $2n$ and
$2m$ (rather than simply $n$ and $m$), respectively, as they are equal to the geometric
distance between the two spins under consideration and the distance from the spin
$\sigma_{(-2m-\frac12,0)}$ to the boundary of the half-plane $\HHd$, respectively.

%% ===========================================================================
%% CROSS-REFERENCES USED BY THE UNIT (as printed; not transcribed here)
%% ---------------------------------------------------------------------------
%% (4.3), (4.4)     : half-plane boundary conditions, Section 4.1
%% Proposition 2.2 / Proposition 2.4 : discrete harmonicity of the Kadanoff-Ceva
%%                    spinors (journal pp. 2286-2289), used for (A.1)-(A.2)
%% Lemma 3.1        : uniqueness of the bounded spinor via optional stopping
%% Theorem A.3 (Wu) : explicit formula (A.4) for D_n, journal p. 2325
%% Lemma A.2        : explicit V via Legendre polynomials, journal p. 2324
%% Theorem 1.1      : the paper's main result, journal p. 2279 (NOT used here:
%%                    the appendix deliberately avoids it, see p. 2323)
%% ===========================================================================

\end{document}
